\section{Discussion}
\label{sec:analysis}

CV predicts connector-level compressibility, but task complexity determines \emph{which} scoring method succeeds within a given connector. On OpenVLA, L2 leads on single-step Object ($+$7.8\,pp at $k{=}0.9$) but the hierarchy inverts on multi-step Long, where Random outperforms both L2 and attention at $k{=}0.7$ (Appendix~\ref{app:extended_discussion}). The monotonic L2 curve versus the non-monotonic random curve on $\pi_{0.5}$ Long (Figure~\ref{fig:cliff_highlight}) reflects position encoding sensitivity: L2 preserves spatial locality after re-indexing, while random creates scattered gaps that maximally disrupt RoPE at $k{\approx}0.9$ before recovering at $k{\leq}0.7$ as surviving positions become approximately uniform (Appendix~\ref{app:extended_discussion}).

\paragraph{Limitations.}
All evaluations use LIBERO simulation; we provide no sim-to-real evidence. The six architectures differ in backbone, encoder, and scale beyond the connector, preventing causal isolation. The CV--SR correlation relies on four closed-loop models, too few for a general law. VTR is validated on SmolVLA only. Most conditions use a single seed; $\pi_{0.5}$ Long uses five-seed means ($n{=}5$, sample s.d.\ ${\leq}$2.7\,pp; Appendix~\ref{app:pi05_long_multiseed}); cross-seed validation on other models shows $\Delta \leq 5.8$\,pp (Appendix~\ref{app:pi05_multiseed},~\ref{app:openvla_multiseed}). Extended discussion in Appendix~\ref{app:extended_discussion}.
